% 3.1 Lighting and energy
\subsection{Lighting and energy}
\label{sec:lighting}

Start with the problem the bottle light only half solves. In sub-Saharan
Africa the share of people with electricity has risen a long way, from
\sspElecAccessStart{} in \sspElecStartYear{} to \sspElecAccess{} in
\sspElecYear{}, and in rural areas it is only \sspRuralElecAccess{}
\parencite{worldbank_wdi}. The standard answer to that gap is the grid:
generate power centrally and run wires to every home. It is a good answer, and
in the long run probably the best one. But it is slow and capital-hungry, and
for a family in a village far from the nearest line, it is not an answer they
can do anything about.

What they could do something about, from around the late 2000s, was a small
solar kit: a panel the size of a book, a battery, a few LED lamps and a USB
port for charging a phone. This is the reframing step in its clearest form.
The question stops being how do we connect this house to the grid?'' and
becomes what does this house need electricity \emph{for}?'' For many
households the first answers are light after dark and a charged phone, and
those need a few watts, not a grid connection. The kits got cheap because the
parts got cheap. Solar cells, LEDs and lithium batteries all fell in price
through the 2010s, for reasons that had nothing to do with African villages
and that I come back to in \cref{sec:against}.

The second reframing was financial. A kit that costs more than a month's
income is out of reach even if it pays for itself in saved kerosene within a
year. Pay-as-you-go companies solved this by letting customers pay in small
instalments, often daily or weekly, through mobile money, with the kit locking
itself if payments stop. One scarcity-born innovation was built on top of
another: without M-Pesa and its cousins (\cref{sec:money}) there would be no
practical way to collect thousands of tiny payments from people without bank
accounts.

The result is real growth. In sub-Saharan Africa the number of people using
off-grid solar grew from about \sspSolarUsersStart{} in \sspSolarStartYear{}
to \sspSolarUsers{} in \sspSolarYear{}, and the number using solar home
systems, the larger kits that come closest to genuine household electricity,
went from \sspShsUsersStart{} to \sspShsUsers{}
\parencite{irena2025offgrid}. Those figures are based mostly on sales reported
by members of the industry association, and therefore undercount. The
whole-market estimate in the official tracking report is that
\sdgOffgridServed{} people worldwide were served by off-grid solar at the
first or second tier of access in \sspSolarYear{}, up from
\sdgOffgridServedPrev{} a year earlier \parencite{esmap2026sdg7}. In
\goglaYear{} alone, industry members sold \goglaKitsSold{} kits
\parencite{gogla2026}.

\begin{figure}[htbp]
\centering
% Figure: electricity access in sub-Saharan Africa, as a share and as a headcount,
% with people using off-grid solar on the same scale.
% Data: data/processed/electricity_access.csv (scripts/process_electricity.py),
%       data/processed/offgrid_solar.csv (scripts/process_offgrid.py).
\begin{tikzpicture}
\begin{groupplot}[
  group style={group size=2 by 1, horizontal sep=0.15\linewidth},
  scarcity axis, year axis,
  width=0.47\linewidth, height=0.42\linewidth,
  xmin=2000, xmax=2024, xtick={2000,2008,2016,2024},
  unbounded coords=jump,
]
  % Left: share of people with access
  \nextgroupplot[
    ymin=0, ymax=100, ytick={0,25,50,75,100},
    title={Share with electricity (\%)},
    title style={font=\sffamily\footnotesize, anchor=south west, at={(0,1.02)}, align=left},
  ]
  \addplot[series context] table[x=year, y=wld, col sep=comma] {../data/processed/electricity_access.csv}
    node[end label, pos=1, text=stone] {World};
  \addplot[series primary] table[x=year, y=ssf, col sep=comma] {../data/processed/electricity_access.csv}
    node[end label, pos=1, text=indigo] {SSA};
  \addplot[series accent] table[x=year, y=ssf_rural, col sep=comma] {../data/processed/electricity_access.csv}
    node[end label, pos=1, text=amber!70!ink, align=left] {SSA\\rural};

  % Right: people without access, and people using off-grid solar (stacked)
  \nextgroupplot[
    ymin=0, ymax=700, ytick={0,200,400,600},
    title={People in SSA (millions)},
    title style={font=\sffamily\footnotesize, anchor=south west, at={(0,1.02)}, align=left},
  ]
  \addplot[series primary] table[x=year, y=ssf_without_m, col sep=comma]
    {../data/processed/electricity_access.csv}
    node[end label, pos=1, anchor=south west, text=indigo] {Without\\electricity};
  \addplot[draw=none, fill=amber!45] table[x=year, y=total, col sep=comma]
    {../data/processed/offgrid_solar.csv} \closedcycle;
  \addplot[draw=none, fill=amber!85!ink] table[x=year, y=shs, col sep=comma]
    {../data/processed/offgrid_solar.csv} \closedcycle;
  \node[end label, text=amber!70!ink, anchor=south west] at (axis cs:2024,82) {Solar\\lights};
  \node[end label, text=amber!70!ink, anchor=south west] at (axis cs:2024,0) {Solar home\\systems};
\end{groupplot}
\end{tikzpicture}

\caption{\takeaway{The share of people with electricity has doubled; the
number without it has not fallen.} In sub-Saharan Africa (SSA), access rose from
\sspElecAccessStart{} in \sspElecStartYear{} to \sspElecAccess{} in
\sspElecYear{}, but population grew almost as fast: about
\sspWithoutElec{} people still had no electricity in \sspElecYear{}, against
\sspWithoutElecStart{} in \sspElecStartYear{}. Headcount is population times
the share without access, rounded to the nearest 5~million; the official
estimate for \sspElecYear{} is \sdgWithoutElecSSA{} \parencite{esmap2026sdg7}.
Shaded bands: people using off-grid solar, which grew from
\sspSolarUsersStart{} in \sspSolarStartYear{} to \sspSolarUsers{} in
\sspSolarYear{}. Solar home systems (11\,W and up, roughly the first rung of
real access) went from \sspShsUsersStart{} to \sspShsUsers{}; the rest are
small lights. These counts rest mostly on sales by members of the industry
association and undercount the whole market, perhaps several times over; some
users are also on the grid and use solar as backup.
\figsource{World Bank, World Development Indicators, series EG.ELC.ACCS.ZS,
EG.ELC.ACCS.RU.ZS and SP.POP.TOTL \parencite{worldbank_wdi}; IRENA,
Off-Grid Renewable Energy Statistics 2025, Africa less North Africa
\parencite{irena2025offgrid}.}}
\label{fig:electricity-access}
\end{figure}
\todo{PAYG and solar-kit impact evidence: Lee, Miguel and Wolfram; Aklin et al.; Grimm et al.}

\begin{counterpoint}
Look at the right-hand panel of \cref{fig:electricity-access} again. The bands
for off-grid solar are growing quickly, but they sit at the bottom of a chart
dominated by the \sspWithoutElec{} people who still have no electricity at all.
The total has barely moved in twenty-five years, because population grew
almost as fast as access. Most of the solar users counted here have small
lights, which fall below even the lowest rung of what the international
tracking framework counts as electricity access, and an estimated
\sdgOffgridBackupShare{} of users are already on the grid and use solar as
backup \parencite{esmap2026sdg7}. Off-grid solar is a real
innovation and a genuine improvement for tens of millions of people. It is not
yet a solution to energy poverty, and it may be closer to a better candle than
to a substitute for the grid.
\end{counterpoint}
